La separación ciega de fuentes es aquella que busca separar las señales originales a partir únicamente de las mezclas observadas, sin información previa sobre las fuentes ni sobre el proceso de mezcla.\cite{cardoso1998blind}

Dos técnicas clásicas son:
\begin{itemize}
    \item \textbf{Análisis de componentes independientes (\textit{ICA} por sus siglas en inglés)} busca componentes estadísticamente independientes, muy usada en aplicaciones de separación de voz.
    
    En aplicaciones de audio, cada señal observada suele corresponder a un micrófono o canal diferente. Por ejemplo, si dos personas hablan simultáneamente y sus voces son registradas mediante dos micrófonos situados en posiciones distintas, cada micrófono recoge una combinación diferente de ambas voces. ICA puede aprovechar estas diferencias para estimar las señales originales.
    Por otro lado, lo más común en el mundo real es que exista solapamiento en cada señal, o suelen compartir el ritmo, la armonía, la tonalidad y la estructura temporal. Por estos motivos el método ICA no suele dar buenos resultados centrándose solamente en buscar elementos estadísticamente independientes y necesita de hipótesis adicionales.

    \item \textbf{\textit{Non-negative Matrix Factorization NMF (Factorización no negativa de matrices)}:} descompone el espectrograma en bases y activaciones.
    
    NMF trata de estimar una matriz de separación y tiene en cuenta el patrón frecuencial asociado, por ejemplo, al timbre de un instrumento, una nota musical o un tipo de evento sonoro. Cada fila de la matriz indica cuándo y con qué intensidad aparece dicho patrón a lo largo del tiempo. 
    Debido a la restricción de no negatividad, la reconstrucción se realiza mediante la suma de contribuciones positivas, sin cancelaciones entre componentes \cite{lee1999learning}. Estas características suelen terminar en mejores resultados.
    Otra diferencia importante es que NMF puede aplicarse a una única mezcla mono. A diferencia de ICA, no necesita obligatoriamente disponer de varias grabaciones de la misma escena sonora. La separación se apoya en la estructura espectral y temporal del audio, y no únicamente en las diferencias espaciales existentes entre varios micrófonos.

\end{itemize}

La versión contraria sería la separación supervisada, donde se aprende teniendo información de los resultados. Es el caso de este trabajo.




